\chapter{Implementation}
\label{ch:implementation}

This chapter describes how the design of Chapter~\ref{ch:design} is realized: the technology stack and repository layout (Section~\ref{sec:stack}), reproducibility engineering (Section~\ref{sec:reprowork}), the three usage surfaces with working examples (Section~\ref{sec:surfaces}), Kaggle ingestion (Section~\ref{sec:kaggle}), quality assurance (Section~\ref{sec:quality}), and the development process itself (Section~\ref{sec:process}).

\section{Technology stack and repository layout}
\label{sec:stack}

The stack was selected for typedness, determinism, and local operation. Python 3.11+ is the implementation language; Pydantic~v2 enforces the contracts; pandas and scikit-learn \cite{pedregosa2011sklearn} provide the data and model layers; the official MCP SDK \cite{mcp2024spec} serves the seven stdio servers; FastAPI and Uvicorn serve the HTTP surface \cite{fastapi2024}; kagglehub ingests public datasets \cite{kagglehub2024}; and SQLite/FTS5 \cite{sqlite2024fts5} backs the context store. The workspace front end is React 18 + TypeScript built with Vite and served as static files by the same FastAPI process, so the entire system runs as a single local service. Table~\ref{tab:stack} lists the pinned core dependencies.

\begin{table}[H]
\centering
\caption{Core dependencies (versions pinned in \texttt{requirements.lock}).}
\label{tab:stack}
\small
\begin{tabular}{llp{7.6cm}}
\toprule
\textbf{Dependency} & \textbf{Role} & \textbf{Why this one} \\
\midrule
pydantic ($\geq$2) & Contracts & Strict typed models with validation; the contract backbone \\
pandas ($\geq$2) & Dataframes & Standard tabular layer \cite{mckinney2010pandas}; profiling and cleaning \\
scikit-learn ($\geq$1.4) & Models & Nine deterministic estimators; \texttt{random\_state} discipline \\
mcp ($\geq$1.6) & Protocol & Official SDK for stdio servers and clients \\
fastapi / uvicorn & HTTP & Async service, SSE streaming for job events \\
kagglehub & Ingestion & First-party Kaggle adapter for public datasets \\
pyarrow & Parquet & Future-proof columnar I/O support \\
matplotlib & Figures & Thesis evidence generation (validity charts, comparison plots) \\
\bottomrule
\end{tabular}
\end{table}

The repository follows a conventional Python layout. The package tree maps one-to-one onto the responsibilities of Table~\ref{tab:modules}; the workspace front end lives in \texttt{web/}; and every runtime output lands under the user's workspace (\texttt{runs/}, \texttt{data/}, \texttt{proposals/}, \texttt{.thelab/}) with persisted references kept relative.

\begin{lstlisting}[language={}, caption={Repository layout (top level, abridged).}, label={lst:layout}, style=cli]
thelab/                  # implementation (10 packages, ~16k LOC)
  contracts/ run/ context/ agents/ ide/ mcp/ sandbox/
  model_service/ eda/ workspace/
  agents/approval.py    # single approval chokepoint
  agents/grounding.py   # shared claim verification
  ide/agentic_round.py  # grounded-autonomy role loop
web/                     # React + TS workspace (Vite build)
tests/                   # 57 test files, 561 test functions
scripts/                 # evaluate_thesis.py, demo_defense.sh, harvest_findings.py,
                         #   probe_model_latency.py, generate_evidence.py, build_ui.sh
examples/                # fixtures, batch configs, Kaggle experiment
docs/                    # user/CLI/API guides, thesis map, legacy records
runs/ data/ proposals/ .thelab/   # runtime workspace (user-local)
thesis/evidence/         # thesis evidence pack (tables, figures, snapshots)
\end{lstlisting}

\section{Reproducibility engineering}
\label{sec:reprowork}

Reproducibility is engineered, not hoped for. Four mechanisms cooperate:

\begin{enumerate}
	\item \textbf{Seed injection.} The seed is a mandatory argument of every run, injected into the estimator (\texttt{random\_state}) and the split; \texttt{build\_estimator} refuses to construct an estimator without it.
	\item \textbf{Manifest provenance.} The manifest records the input dataset hash, the full training configuration, the seed, and the resolved dependency versions (Python, pydantic, pandas, scikit-learn, numpy, joblib) --- captured automatically at run time, not maintained by hand.
	\item \textbf{Lockfile.} \texttt{requirements.lock} pins the environment; the documented install path builds the virtual environment from the lock, making the recorded dependency versions reproducible in practice (with the documented caveat that the lock is not hash-verified; Section~\ref{sec:limitations}).
	\item \textbf{Generated notebooks.} Every completed run can produce a reproducible research notebook artifact (\texttt{report.ipynb}, generated on demand over \texttt{GET /runs/\{id\}/notebook}): title and manifest summary, dataset load cell, a reproduce cell invoking the pipeline with the exact recorded parameters, metrics and validation display, artifact index, and findings from sub-agent interpretations when present. The notebook is built as standards-compliant nbformat JSON without new dependencies, and is viewable read-only in the UI.
\end{enumerate}

\section{Three ways in: CLI, Python API, HTTP}
\label{sec:surfaces}

The same pipeline is reachable three ways. The CLI (Listing~\ref{lst:cli}) is the primary surface for scripting and the reference for documentation.

\begin{lstlisting}[language=bash, caption={Reproducible training and prediction from the CLI.}, label={lst:cli}, style=cli]
# Train one model deterministically
thelab run model --dataset examples/iris.csv --target species \
  --model logistic_regression --seed 42 --output runs

# Output:
# Run completed: run-20260814-222154-a5a01046
#   Model: logistic_regression   Seed: 42
#   Test accuracy: 1.000000      Test macro F1: 1.000000

# One-off prediction from an approved run
thelab predict --run-id <run_id> --features '5.1,3.5,1.4,0.2'

# Index and search the local evidence store
thelab context index --source .thelab/local-logs/agent-events.jsonl
thelab context search "reproducibility" --limit 10

# Agent CLI: worker mode proposes an experiment (EDA-grounded)
thelab-agent --mode worker --dataset examples/iris.csv --target species \
  --model-grid logistic_regression,random_forest --seeds 42,43 \
  --provider openrouter
\end{lstlisting}

The Python API offers the same capabilities to notebooks and scripts (Listing~\ref{lst:python}): a low-level \texttt{run\_model} and \texttt{try\_all\_models}, and a high-level \texttt{thelab.quick} facade.

\begin{lstlisting}[language=python, caption={Python API: deterministic run, try-all comparison, and a proposal from a notebook.}, label={lst:python}, style=estilo]
from thelab.quick import experiment, compare

exp = experiment("examples/iris.csv", target="species", model="svc")
print(exp.metrics)                       # recorded metrics
print(exp.predict([5.1, 3.5, 1.4, 0.2])) # inference from the run

for e in compare("examples/iris.csv", target="species"):
    print(e.model, e.metrics.get("test_accuracy"))

# Agent proposals are importable too -- deterministic fallback, no LLM needed
import asyncio
from thelab.agents.mock import MockProvider
from thelab.agents.worker import WorkerAgent
from thelab.ide.proposals_api import approve_and_run_proposal

worker = WorkerAgent(provider=MockProvider([]), servers=[],
                     proposals_dir="proposals")
p = asyncio.run(worker.propose(goal="Predict species",
                               dataset="examples/iris.csv",
                               target="species",
                               model_grid=["logistic_regression"],
                               seeds=[42]))
outcome = approve_and_run_proposal(p.proposal_id, principal="notebook")
\end{lstlisting}

The HTTP surface serves the UI and a JSON API with a uniform \texttt{\{ok, data\}} envelope (Listing~\ref{lst:http}); long operations run as background jobs with SSE progress and cooperative cancellation.

\begin{lstlisting}[language=bash, caption={HTTP surface: training and inference via the local service.}, label={lst:http}, style=cli]
thelab-model-service --host 127.0.0.1 --port 8000

curl -s -X POST localhost:8000/train -H 'Content-Type: application/json' \
  -d '{"dataset_id": "fixtures/iris.csv", "target": "species",
       "model": "logistic_regression", "seed": 42}'

curl -X POST http://127.0.0.1:8000/predict -H "Content-Type: application/json" \
  -d '{"run_id": "run-20260829-...", "features":
       [{"sepal_length": 5.1, "sepal_width": 3.5,
         "petal_length": 1.4, "petal_width": 0.2}]}'
\end{lstlisting}

Independently of these three, an MCP client consumes the factory through the servers of Table~\ref{tab:mcpservers}; the bundled demo client smoke-tests the round trip without any MCP SDK knowledge beyond the protocol itself:

\begin{lstlisting}[language=bash, caption={MCP smoke tests over stdio.}, label={lst:mcpdemo}, style=cli]
thelab-mcp-demo model_registry --run-id <run_id>  # list_models, card, predict
thelab-mcp-demo data_catalog  --run-id <run_id>   # datasets, profiles, contracts
thelab-mcp-demo context                           # search the evidence store
\end{lstlisting}

The positioning rule of thumb documented in the user guide: \emph{use the UI to explore, the CLI to automate, MCP when another agent needs the factory}.

\section{Kaggle ingestion and context packs}
\label{sec:kaggle}

Public data is a first-class input. From a Kaggle slug, the ingestion path (Listing~\ref{lst:kagglecli}) downloads the dataset into \texttt{data/uploads/} through kagglehub (the only networked step, always explicit user action), fetches the dataset's public page context (description, tags, file metadata --- parsed from the page's embedded state, failing soft), merges it with a local profile (shape, dtypes, head), and persists the merged \emph{context pack} alongside the CSV (\texttt{<name>.kaggle.json}). The pack is indexed into the context store as a searchable event and exposed to the chat agent through the \texttt{get\_dataset\_context} tool, so proposals are grounded in the dataset's own documentation --- for example, ordinal encodings documented on the Kaggle page become part of the agent's evidence.

\begin{lstlisting}[language=bash, caption={Kaggle ingestion via CLI.}, label={lst:kagglecli}, style=cli]
thelab ingest kaggle pavansubhasht/ibm-hr-analytics-attrition-dataset
# -> downloads CSV to data/uploads/, builds + indexes the context pack
\end{lstlisting}

\section{Testing and quality assurance}
\label{sec:quality}

Quality gates are enforced by four automated layers, run before any change is considered complete:

\begin{lstlisting}[language=bash, caption={The four quality gates.}, label={lst:gates}, style=cli]
.venv/bin/ruff check thelab tests scripts        # lint (E,F,I,UP,W,B)
.venv/bin/mypy thelab                            # static types
.venv/bin/python -m pytest tests/ -q             # 495 tests
PATH=.venv/bin:$PATH python scripts/evaluate_thesis.py   # RQ1-RQ3
\end{lstlisting}

The test suite covers every package at unit level and exercises the full stack end-to-end against the mock provider --- including a dedicated real-data-hardening suite that replays the S\&P-scale scenarios, and a suite that tests the thesis evaluator itself. Notable integration tests verify the approval boundary end-to-end (proposal $\rightarrow$ approve $\rightarrow$ batch $\rightarrow$ runs), the sandbox policy against adversarial inputs, the context store's immutability under reads, and the MCP contract of all seven servers.

\section{Development process}
\label{sec:process}

The Lab was developed with heavy AI assistance under a process designed to keep the author in control --- a decision with its own history. An initial attempt to generate the first implementation in one automated pass (a code agent plus an audit agent, at a cost of approximately EUR~30) produced code the author could not read or confidently own. The process was redesigned around a rule the author summarized as \emph{slice implemented -- slice read}: every implemented slice must be read and understood before the next begins. The project's subsequent working mode --- small, focused slices; contracts before features; four quality gates before done --- is the direct descendant of that lesson, and the development log records its effect on cost structure as well: audit-style agent runs cost \$0.34--\$4 per run, while large-context coding-agent runs reached $\sim$\$10 each, an asymmetry that motivated the review-heavy, generation-light workflow.

The agentic tooling itself was part of the stack under test: coding agents implemented slices against the typed contracts, audit agents (a role rotated across models after one frontier model repeatedly stalled) reviewed diffs adversarially, and a read-only external assistant served as a second opinion. This thesis therefore reports not only a system built for agentic ML, but evidence gathered from building a system \emph{with} agentic tooling under reproducible process constraints.

Three decisions from the pre-implementation planning phase survive, unchanged, in the delivered system --- recorded here as provenance of the design:

\begin{itemize}
	\item \textbf{The lab metaphor.} The project was framed from the first week as ``infrastructure is the lab; the custom apps are the experiments you run in it'' --- the direct origin of the system's name and of the decision to build the deterministic core as a reusable factory rather than a single-purpose application.
	\item \textbf{The descope checkpoint.} The phased plan set an explicit trip-wire: if the core loop (MCP servers, pipeline, evaluator) was not stable by early August, the semantic-memory layer would be descoped to ``designed, partially implemented'' rather than risked --- a decision exercised exactly as planned, and the reason the context store is keyword-only (FTS5) today (Section~\ref{sec:bg-context}).
	\item \textbf{The deterministic judge.} Planning-phase sketches of a ``deterministic judge'' for agent outputs --- validation with parameters, not LLM self-assessment --- crystallized into the dual design of the deterministic pipeline (which judges runs) and the grounding checks of the harness (which judge agent claims).
\end{itemize}

A first-review audit of the near-complete system (September 2026) extended this discipline inward. It found the engineering core solid but the agentic claims partially overstated: sub-agents shared one system prompt with an unused role argument, some orchestration paths self-approved their own proposals (bypassing the intended human boundary), and threaded user feedback was never read. Every finding was closed before any new capability was built: per-role prompt contracts replaced the shared prompt, a single approval chokepoint made agent-initiated execution impossible without explicit human approval (Section~\ref{sec:approvalgate}), feedback became a real consumer in all three stage interpretations, and the dead code was deleted. The episode is recorded here deliberately: in agentic systems, the gap between what the code does and what the documentation claims is itself a failure mode, and closing it --- 495 tests green, deterministic path untouched --- preceded every new agentic claim in this thesis.
